\subsection{Checklist: Code of Ethics}
\label{app:checklist_code_of_ethics}

\paragraph{Question.}
Does the research conducted in the paper conform, in every respect, with the
NeurIPS Code of Ethics?

\paragraph{Answer (Yes/No/NA).}
Yes.

\paragraph{Justification.}
The paper studies sparse Koopman autoencoders for forecasting and interpreting
nonlinear dynamical systems. The experiments reported in the manuscript use
simulation and benchmark trajectory data: a procedurally generated collection
of two-dimensional multibasin dynamical systems, and a ten-system
three-dimensional Dysts forecasting benchmark from the pinned project
environment. The models observe stored trajectory states. They do not collect,
process, infer from, or release human-subject data, personal information,
private records, medical data, demographic attributes, user-generated content,
or sensitive social data.

The manuscript also separates training-time assumptions from benchmark-only
evaluation information. In the intended training and deployment setting, the
method is not given basin labels, the number of basins, or trajectory-to-basin
assignments. When benchmark basin labels or attractor metadata are available,
they are used after training for evaluation slices, diagnostic scoring,
controlled-transfer construction, and architecture diagnostics, not as ordinary
training supervision for the support-selection mechanism. This keeps the paper's
claims focused on label-free representation learning from simulated dynamical
trajectories rather than on decisions about people or protected groups.

\paragraph{Anonymization and ethics considerations.}
The manuscript is prepared under anonymous authorship. Because the experiments
use synthetic or public benchmark dynamical-system trajectories rather than
participant data, no participant de-identification, consent process, or
institutional human-subject review is implicated by the reported experiments.
The procedural benchmark construction and the Dysts source should remain
clearly disclosed and cited, including the package version used for the
paper-facing benchmark. Any released code, logs, tables, or supplementary
artifacts should be checked before submission for accidental disclosure of
author identity, local filesystem paths, usernames, cluster job metadata, or
other non-anonymous provenance.

\paragraph{Caveats.}
The paper is methodological and simulation-based. Its forecasting and
interpretability results should not be presented as validated evidence for
safety-critical control, clinical decision support, surveillance, or other
deployment contexts where prediction errors could cause direct harm. If future
versions add real-world datasets, human-generated measurements, proprietary
data, biological or medical data, or safety-critical deployment claims, the
ethics analysis should be revisited for consent, privacy, licensing,
dual-use risk, and domain-specific review requirements. If code or artifacts
are released, the authors should also verify external benchmark licensing and
venue disclosure requirements for generated benchmark content and compute use.
